
We opened with a counterintuitive finding: spurious-minority samples appear \textit{more} aligned with the within-batch mean gradient than majority samples on Waterbirds at epoch 1 (alignment ROC-AUC = 0.150), despite the theoretical prediction that conflicting gradient directions should reveal them. This inversion is consistent with batch contamination by high-magnitude minority gradients.

The results are clear: within-batch gradient alignment ROC-AUC $\leq$ 0.35 (Waterbirds) and $\leq$ 0.63 (CelebA) against per-sample loss 0.77--0.97 under identical conditions. The failure is not marginal. Per-sample loss remains the practical baseline for annotation-free minority identification.

The inversion is a diagnosis, not a dead end. The within-batch mean is contaminated by high-loss minority samples --- structurally, due to mini-batch imbalance. The natural corrective direction is a two-pass global mean gradient, stable and majority-dominated, requiring one additional full-dataset pass per epoch, annotation-free. Distinguishing batch contamination from last-layer gradient compression (the primary alternative explanation) requires penultimate-layer alignment experiments --- an immediate next step.

\textbf{Established:} within-batch alignment fails, with a novel inversion failure mode; the failure is consistent with a mechanistic and prevalence-dependent contamination effect; two-pass global mean is the expected corrective direction. \textbf{Open:} whether the correction resolves the inversion; whether penultimate-layer alignment provides complementary information. Gradient direction as a minority signal remains theoretically sound. The reference direction was the problem --- not the concept.

